\section{仿真验证、任务完成与时间最优}

\label{sec:verify}

全部结果由本地模拟器（\texttt{jammers-py}，与官方引擎同判据的离线复现）在
\texttt{problem\_no=4} 场景下、\texttt{T4.py --practice 20 --seed 0} 产生：20 局、
种子 $0$--$19$，每局场景布局与示向度噪声完全确定、逐字节可复现。

\subsection{任务完成结果（20 局全量）}

\begin{table}[H]
  \centering
  \caption{20 局演练汇总（seed 0--19，20 点定案布局）}
  \label{tab:survey}
  \begin{tabular}{lcc}
    \toprule
    \textbf{指标} & \textbf{数值} & \textbf{备注} \\
    \midrule
    清除比例 & $\mathbf{256/256 = 100\%}$ & 20/20 局全部清完，0 失败 \\
    平均虚拟总耗时 & $6445$ s & $\approx 107$ min，逐局 $5580$--$7861$ s \\
    平均里程 & $24\,392$ m & 其中检测扫描 $16\,996$ m \\
    平均测向 & $247$ 次/局 & 顺路清除尝试 $8.6$ 个（命中 $7.1$ + 未中 $1.5$） \\
    定位误差 & 均值 $7.66$ m & 最差单源 $19.67$ m，全部 $\le 20$ m \\
    首次听到步骤 & $\le 19$ & 共 20 个测量位置，最晚第 19 步 \\
    \bottomrule
  \end{tabular}
\end{table}

\begin{figure}[H]
  \centering
  \includegraphics[width=0.78\textwidth]{figures/fig_t4_episode_clear.pdf}
  \caption{20 局清除数（$256/256$）与各局虚拟时间。}
  \label{fig:clear}
\end{figure}

\begin{figure}[H]
  \centering
  \includegraphics[width=0.78\textwidth]{figures/fig_t4_first_heard.pdf}
  \caption{听到率实证：全部 $256$ 个源在扫描第 $\le 19$ 步内首次听到。}
  \label{fig:firstheard}
\end{figure}

\begin{figure}[H]
  \centering
  \includegraphics[width=0.78\textwidth]{figures/fig_t4_localization_error.pdf}
  \caption{清除点定位误差：局均 $7.66$ m，全部落在 $20$ m 清除半径内。}
  \label{fig:locerr}
\end{figure}

逐局明细（seed : 清除数/源数（虚拟秒））：

\begin{quote}\small
\texttt{seed0=16/16(6238s)  seed1=11/11(6544s)  seed2=16/16(6585s)  seed3=11/11(6246s)\\
seed4=11/11(5648s)  seed5=14/14(6064s)  seed6=16/16(7166s)  seed7=12/12(7080s)\\
seed8=11/11(6217s)  seed9=13/13(7861s)  seed10=14/14(5966s)  seed11=13/13(6212s)\\
seed12=13/13(6582s)  seed13=12/12(6486s)  seed14=10/10(6191s)  seed15=11/11(6027s)\\
seed16=12/12(5580s)  seed17=14/14(6209s)  seed18=11/11(7133s)  seed19=15/15(6859s)}
\end{quote}

\subsection{"方向不对"的实证}

seed 0 频道 6：该源在第 $0$--$19$ 个测量位置的多次测量中全部 \texttt{no\_signal}（源真实
存在、距离未必超限，只是机器狗始终站在波束背面），直到扫描后期从正确方向命中才首次听到，
随后一次清除成功。若没有 12 个外圈测量点，这个源会被当成"不存在"而漏掉——这正是问题四
检测布局的意义。

\subsection{时间构成与 6445 s 下界}

虚拟时间构成（图~\ref{fig:cost}）：移动 $\sim 78\%$（其中检测扫描 $17.0$ km 为检测必须、
约 $54\%$）、测向 $\sim 17\%$（大半是"每个测量位置测未听频道"的必需试探）、清除
$\sim 5\%$。前两项合计约 $90\%$ 本质不可压缩。

\begin{figure}[H]
  \centering
  \includegraphics[width=0.78\textwidth]{figures/fig_t4_cost.pdf}
  \caption{时间构成：17.0 km 检测扫描行驶为固定搜索成本。}
  \label{fig:cost}
\end{figure}

在此约束下逐个尝试可调旋钮（2--20 局同 seed 对拍，表~\ref{tab:knobs}），收敛到
$6445$ s 的实测时间下界。

\begin{table}[H]
  \centering
  \caption{时间优化旋钮的系统探索（全部经 2--20 局同 seed 对拍裁决）}
  \label{tab:knobs}
  \begin{tabular}{lll}
    \toprule
    \textbf{旋钮} & \textbf{尝试} & \textbf{实测结论} \\
    \midrule
    途中顺路清除 开/关 & 开关对比 & \textbf{开更省}（2 局省 150--200 s）：命中省收尾专程 $>$ 绕路，保留 \\
    顺路排除半径 600 m & 600 $\to$ 0 m & 0 时内部源偏离外圈路线、绕路 $>$ 命中收益，20 局 +178 s，回滚 \\
    收尾阶段顺路清除 & 途中试清其余 & seed0 持平、seed1 +339 s：收尾剩的多是顽固源，未中 3 s 白花，否 \\
    扫描顺序"外圈优先" & 先外圈后内部 & 路线 +1.4 km、顺路命中骤降、时间 +1240--2893 s，否 \\
    顺路清除横向绕路上限 & 只清贴近直线点 & seed 间方向相反（$-322$ s / $+256$ s），固定阈值不稳，不设 \\
    OBS\_CAP 4 $\to$ 3 & 每源少采 1 条示向度 & 测向仅 $-5$ 次/局，有定位精度风险，保持 4 \\
    \bottomrule
  \end{tabular}
\end{table}

\textbf{结论：$6445$ s $\approx$ 时间下界的实现值}（20 点定案；曾定案的 23 点版 $6688$ s，
多 3 个内部补点的测向代价。题目时限约束的是现实时钟——测试窗口 25 min / 程序 20 min——
与虚拟 100 h 上限；模拟器即时演算，107 min 虚拟时间在现实秒级内跑完，不构成超时）。

\subsection{神经网络独立交叉验证}

为排除"人工布局不是最优"的疑虑，用神经网络代理 + 遗传算法独立搜索测量点布局（
\texttt{cumcm/t4/nn\_layout\_speed.py}）：固定 7 覆盖基点，让外圈点个数与摆位自由进化，
精英精评 = $40$ 万蒙特卡洛听率 + 贴边对抗 + TSP 路线；最后全部用 20 局演练裁决。

\begin{figure}[H]
  \centering
  \includegraphics[width=0.95\textwidth]{figures/fig_t4_nn_layout_compare.png}
  \caption{神经网络搜索布局 vs 人工定案布局对比（左侧人工 20 点；右侧神经网络自由搜索
  布局——外圈点不均匀、路线更长）。}
  \label{fig:nn_compare}
\end{figure}

\begin{figure}[H]
  \centering
  \includegraphics[width=0.95\textwidth]{figures/fig_t4_nn_speed_compare.png}
  \caption{各布局 20 局实测平均虚拟时间：人工 20 点版 $6445$ s 为最小，神经网络全部尝试
  （自由、12 外圈、11 外圈）均更慢。}
  \label{fig:nn_speed}
\end{figure}

\begin{table}[H]
  \centering
  \caption{20 局全量裁决：神经网络全部搜索方向均无法击败人工 20 点版}
  \label{tab:nn_verdict}
  \begin{tabular}{lccc}
    \toprule
    \textbf{布局} & \textbf{位置数} & \textbf{听率} & \textbf{20 局平均虚拟时间} \\
    \midrule
    人工 23 点（历史候选）          & 23 & $99.9974\%$ & $6688$ s \\
    \textbf{人工 20 点（定案）}     & 20 & $99.9891\%$ & $\mathbf{6445}$ s（$256/256$）\\
    神经网络 12 外圈（非均匀摆位）  & 20 & $99.9218\%$ & $6947$ s（$256/256$）\\
    神经网络 11 外圈（非均匀摆位）  & 19 & $99.7778\%$ & $6724$ s（$\mathbf{255/256}$，漏 1 源）\\
    神经网络自由 23 位置            & 23 & $99.9419\%$ & $7697$ s（2 局：6960 / 8433 s）\\
    \bottomrule
  \end{tabular}
\end{table}

\textbf{结论}：神经网络独立搜索证伪了"存在更短耗时布局"的假设——听率是隐性时间成本：听率从
$99.9891\%$ 降到 $99.7778\%$（11 外圈）时，20 局演练里\textbf{真的有 1 个源始终没被听到、$255/256$
漏清}；NN 布局实测总里程 $25.4$--$29.6$ km 也高于人工的 $24.4$ km。可见均匀对称外圈已是顺路清除
路径的实战极值。这使 $6445$ s 的下界可信，也最终确立了第~\ref{sec:detect} 节的 20 点定案布局。